\documentclass[10pt,a4paper]{amsart}
\usepackage[margin=19mm]{geometry}
\usepackage{amsmath,amssymb,mathtools}
\usepackage[scr=rsfs,cal=euler]{mathalfa}
\usepackage{xfrac}
\usepackage[unicode,hidelinks,bookmarks=false]{hyperref}
\newcommand{\ie}{i.e.\ }
\newcommand{\cf}{cf.\ }
\newcommand{\resp}{respectively\ }
\newcommand{\Subsection}{\subsection}
\providecommand{\nobreakdash}{\nobreak\hbox{-}}
\setlength{\emergencystretch}{2em}
\multlinegap0pt
\usepackage{amssymb}
\usepackage{xcolor}
\usepackage{mdframed}
\usepackage{tcolorbox}
\newtheorem{theorem}{Theorem}
\newcommand{\Ddim}{\mathbb{R}^D}
\title{\large \textbf{Identifiability of the unnormalized graph Laplace operators}}
\author{Susovan Pal}
\date{Source: arXiv:2601.19589v1 (CC BY 4.0)}
\begin{document}
\maketitle

\noindent\emph{Source and licence.} \large \textbf{Identifiability of the unnormalized graph Laplace operators}. Version arXiv:2601.19589v1 (CC BY 4.0), distributed by arXiv under the Creative Commons Attribution 4.0 International licence, http://creativecommons.org/licenses/by/4.0/, as recorded in the arXiv licence metadata for this exact version and shown on the abstract page https://arxiv.org/abs/2601.19589v1. Full licence notice: \texttt{licenses/arxiv-2601.19589-CC-BY-4.0.txt}.\par\smallskip

\noindent\textit{Editorial scope.} Counted unit: the article's theorem thm:extrinsic\_induced\_metric\_identifiable (source0.tex lines 663--694). The article's complete original proof of it is delivered with it (source0.tex lines 696--779, one \texttt{proof} environment of 84 lines). No mathematics is written, repaired or re-derived: the statement and the proof are the article's own lines, in the article's own notation, and no retained line is substituted or re-flowed.  Retained uncounted as the source-local context the proof needs: lines 360--375.  Nothing was omitted from any retained span, so the package carries no omission record.  No retained line contains a citation, so the wrapper carries no bibliography and no bibliography entry.  No retained line contains a figure environment, an image-inclusion command or drawing code, so no image is delivered and the package declares no asset dependency.  Editorial formatting only, in addition: an AMS \texttt{amsart} preamble that restates the article's own declarations and macros used by the retained lines, namely the environments \texttt{theorem} on separate counters, together with the standard packages those lines need.  The retained environments print in this document as Theorem 1 in order of appearance, whereas the article prints the same items with the numbering of its own full document.  The complete native source of the article as distributed in the arXiv e-print for this exact version is bundled unchanged as \texttt{sources/193466/source0.tex}.
\begin{theorem}[\textbf{Intrinsic operator identifies the Riemannian metric (for fixed bandwidth and positive density)}]
\label{thm:intrinsic_metric_identifiable}
Let $M$ be a smooth compact manifold without boundary, $\dim M=d$.
Fix $t>0$ and $p\in C(M)$ with $p(x)>0$ for all $x\in M$.
For any $C^2$ Riemannian metric $g$ on $M$, define
\[
(L^{(g)}_{t,p} f)(x)
:=\frac{1}{t^{\frac d2+1}}\int_M \exp\!\Big(-\frac{d_g(x,y)^2}{t}\Big)\,\bigl(f(x)-f(y)\bigr)\,p(y)\,d\mu_g(y),
\qquad f\in C(M).
\]
If $g_1,g_2$ are $C^2$ metrics on $M$ such that
\[
L^{(g_1)}_{t,p}=L^{(g_2)}_{t,p}\quad\text{as operators on }C(M),
\]
then $g_1=g_2$.
\end{theorem}

\begin{theorem}[\textbf{Extrinsic operator identifies the embedding-induced Riemannian metric for fixed bandwidth}]
\label{thm:extrinsic_induced_metric_identifiable}
Let $M$ be a compact smooth manifold without boundary, $\dim M=d$.
Fix $t>0$ and a continuous density $p(x)>0$ for all $x\in M$.
Let $\iota_1,\iota_2:M\to {\Ddim}$ be $C^2$ embeddings and denote the induced metrics
\[
g_i:=\iota_i^*\langle\cdot,\cdot\rangle_{{\Ddim}},
\qquad \mu_i:=\mu_{g_i},
\qquad 
K_i(x,y):=\exp\!\Big(-\frac{\|\iota_i(x)-\iota_i(y)\|^2}{t}\Big),
\quad i=1,2.
\]
Define the corresponding extrinsic continuous graph Laplace operators by
\[
(L^{\mathrm{ext}}_{t,p}(g_i;\iota_i)f)(x)
:=\frac{1}{t^{\frac d2+1}}
\int_M K_i(x,y)\,\bigl(f(x)-f(y)\bigr)\,p(y)\,d\mu_i(y),
\qquad f\in C(M).
\]
If
\[
L^{\mathrm{ext}}_{t,p}(g_1;\iota_1)=L^{\mathrm{ext}}_{t,p}(g_2;\iota_2)
\quad\text{as operators on }C(M),
\]
then
\[
\|\iota_1(x)-\iota_1(y)\|
=\|\iota_2(x)-\iota_2(y)\|
\quad\text{for all }x,y\in M,
\]
and consequently $g_1=g_2$.
\end{theorem}

\begin{proof}
For $i=1,2$ define the integral operators
\[
(T_i f)(x):=\int_M K_i(x,y)\,f(y)\,p(y)\,d\mu_i(y),
\qquad f\in C(M).
\]
Then
\[
(L^{\mathrm{ext}}_{t,p}(g_i;\iota_i)f)(x)
=\frac{1}{t^{\frac d2+1}}
\Big((T_i\mathbf 1)(x)\,f(x)-(T_i f)(x)\Big).
\]
Hence the assumption
$L^{\mathrm{ext}}_{t,p}(g_1;\iota_1)=L^{\mathrm{ext}}_{t,p}(g_2;\iota_2)$
implies that for all $f\in C(M)$ and all $x\in M$,
\begin{equation}
\label{eq:ext_mult_identity_induced_rewrite}
(T_1 f)(x)-(T_2 f)(x)
=\big((T_1\mathbf 1)(x)-(T_2\mathbf 1)(x)\big)\,f(x).
\end{equation}

\smallskip
\noindent\textbf{Step 1: reduction to equality of the integral operators.}
The argument showing from \eqref{eq:ext_mult_identity_induced_rewrite} that
\[
T_1\mathbf 1=T_2\mathbf 1
\quad\text{and hence}\quad
T_1=T_2
\]
is \emph{identical} to \textbf{Step~1} in the proof of
Theorem~\ref{thm:intrinsic_metric_identifiable}
(the Riesz representation / Dirac mass argument using non-atomicity of the
Riemannian volume measures), and is therefore omitted.

\smallskip
\noindent\textbf{Step 2: identification of the distance in the ambient space.}
Fix $x\in M$. Since $T_1=T_2$, for all $f\in C(M)$,
\[
\int_M K_1(x,y) f(y)\,p(y)\,d\mu_1(y)
=
\int_M K_2(x,y) f(y)\,p(y)\,d\mu_2(y).
\]
Because $g_1$ and $g_2$ are smooth Riemannian metrics on compact $M$,
their volume measures are mutually absolutely continuous; write
$d\mu_2=w\,d\mu_1$ with a continuous function $w>0$.
Then
\[
\int_M f(y)\,\bigl(K_1(x,y)-K_2(x,y)\,w(y)\bigr)\,p(y)\,d\mu_1(y)=0
\qquad(\forall f\in C(M)).
\]
Since $p>0$ and $\mu_1$ has full support, this implies, using continuity of $K_i, w, i=1,2$
\[
K_1(x,y)=K_2(x,y)\,w(y)\qquad\text{for all }y\in M.
\]
Setting $y=x$ yields $1=1\cdot w(x)$, hence $w(x)=1$ for all $x$ and thus
$w\equiv 1$. Consequently $K_1\equiv K_2$ on $M\times M$, and this gives 
\(
\|\iota_1(x)-\iota_1(y)\|^2
=\|\iota_2(x)-\iota_2(y)\|^2
\text{ for all }x,y\in M.
\)

\smallskip
\noindent\textbf{Step 3: identification of the induced metric.}
Fix $x\in M$ and $v\in T_xM$. Let $\gamma:(-\varepsilon,\varepsilon)\to M$
be a $C^2$ curve with $\gamma(0)=x$ and $\dot\gamma(0)=v$.
For $i=1,2$, Taylor expansion yields
\[
\iota_i(\gamma(s))
=\iota_i(x)+s\,d\iota_i(x)[v]+o(s)\qquad(s\to 0),
\]
and therefore
\[
\|\iota_i(\gamma(s))-\iota_i(x)\|^2
=s^2\,\|d\iota_i(x)[v]\|^2+o(s^2)\qquad(s\to 0).
\]
By Step~2 the left-hand sides coincide for $i=1,2$ for all sufficiently small
$s$, hence
\[
\|d\iota_1(x)[v]\|^2=\|d\iota_2(x)[v]\|^2.
\]
Since $g_i(x)(v,v)=\|d\iota_i(x)[v]\|^2$, this holds for all $x$ and $v$,
and we conclude that $g_1=g_2$.
\end{proof}

\end{document}
